
\documentclass{article}
\usepackage{colortbl}
\usepackage{makecell}
\usepackage{multirow}
\usepackage{supertabular}

\begin{document}

\newcounter{utterance}

\centering \large Interaction Transcript for game `dond', experiment `semi\_it', episode 14 with claude{-}sonnet{-}4{-}20250514{-}t1.0.
\vspace{24pt}

{ \footnotesize  \setcounter{utterance}{1}
\setlength{\tabcolsep}{0pt}
\begin{supertabular}{c@{$\;$}|p{.15\linewidth}@{}p{.15\linewidth}p{.15\linewidth}p{.15\linewidth}p{.15\linewidth}p{.15\linewidth}}
    \# & $\;$A & \multicolumn{4}{c}{Game Master} & $\;\:$B\\
    \hline

    \theutterance \stepcounter{utterance}  
    & & \multicolumn{4}{p{0.6\linewidth}}{
        \cellcolor[rgb]{0.9,0.9,0.9}{
            \makecell[{{p{\linewidth}}}]{
                \texttt{\tiny{[P1$\langle$GM]}}
                \texttt{State giocando a un gioco di negoziazione in cui dovete accordarvi su come dividere una serie di oggetti tra voi e un altro giocatore.} \\
\\ 
\texttt{Regole:} \\
\texttt{(a) A Lei e all'altro giocatore viene dato un insieme di oggetti. Ognuno di voi riceve anche una funzione di valore segreta, che rappresenta il valore di ciascun tipo di oggetto.} \\
\texttt{(b) Si scambiano messaggi con l'altro giocatore per concordare chi si aggiudica gli oggetti. Potete inviare un massimo di 5 messaggi ciascuno, oppure terminare in anticipo facendo una proposta segreta in qualsiasi momento.} \\
\texttt{(c) A ciascuno di voi viene chiesto di inviare una proposta segreta indicando gli oggetti che desiderate, formattata tra parentesi quadre come segue: "{[}Proposta: <numero> <nome oggetto>, <numero> <nome oggetto>, <...>{]}".} \\
\texttt{(d) Se le vostre proposte sono complementari, cioè ci sono abbastanza oggetti per soddisfare entrambe le proposte, a ciascun giocatore viene assegnato un punteggio basato sulla somma dei valori degli oggetti ricevuti. In caso contrario, entrambi ricevono zero punti.} \\
\texttt{(e) L'obiettivo è massimizzare il proprio punteggio.} \\
\\ 
\texttt{Cominciamo.} \\
\\ 
\texttt{L'insieme degli oggetti disponibili è:} \\
\\ 
\texttt{1 cornice, 1 scarpa, 3 martelli} \\
\\ 
\texttt{La funzione valore segreta è:} \\
\\ 
\texttt{cornice: 8, scarpa: 2, martello: 0} \\
\\ 
\texttt{IMPORTANTE: i vostri messaggi, a meno che non si tratti di una proposta segreta, vengono trasmessi direttamente all'altro giocatore, quindi non includete alcuna risposta alle regole o testo di annuncio del vostro messaggio. Per fare una proposta segreta, utilizzate il formato indicato. Non utilizzare delle parentesi quadre quando si comunica all'altro giocatore, altrimenti verrà interpretata come una proposta segreta.} \\
            }
        }
    }
    & \\ \\

    \theutterance \stepcounter{utterance}  
    & \multicolumn{4}{p{0.6\linewidth}}{
        \cellcolor[rgb]{0.9,0.9,0.9}{
            \makecell[{{p{\linewidth}}}]{
                \texttt{\tiny{[P1$\rangle$GM]}}
                \texttt{Ciao! Vedo che abbiamo alcuni oggetti interessanti da dividere. Che ne pensi se iniziamo a discutere le nostre preferenze? Io sono particolarmente interessato alla cornice, ma sono aperto a negoziare su tutto. Tu hai qualche preferenza particolare tra questi oggetti?} \\
            }
        }
    }
    & & \\ \\

    \theutterance \stepcounter{utterance}  
    & & & \multicolumn{2}{p{0.3\linewidth}}{
        \cellcolor[rgb]{0.9,0.9,0.9}{
            \makecell[{{p{\linewidth}}}]{
                \texttt{\tiny{[GM$|$GM]}}
                \texttt{continue} \\
            }
        }
    }
    & & \\ \\

    \theutterance \stepcounter{utterance}  
    & & & \multicolumn{2}{p{0.3\linewidth}}{
        \cellcolor[rgb]{0.9,0.9,0.9}{
            \makecell[{{p{\linewidth}}}]{
                \texttt{\tiny{[GM$|$GM]}}
                \texttt{Ciao! Vedo che abbiamo alcuni oggetti interessanti da dividere. Che ne pensi se iniziamo a discutere le nostre preferenze? Io sono particolarmente interessato alla cornice, ma sono aperto a negoziare su tutto. Tu hai qualche preferenza particolare tra questi oggetti?} \\
            }
        }
    }
    & & \\ \\

    \theutterance \stepcounter{utterance}  
    & & \multicolumn{4}{p{0.6\linewidth}}{
        \cellcolor[rgb]{0.9,0.9,0.9}{
            \makecell[{{p{\linewidth}}}]{
                \texttt{\tiny{[GM$\rangle$P2]}}
                \texttt{State giocando a un gioco di negoziazione in cui dovete accordarvi su come dividere una serie di oggetti tra voi e un altro giocatore.} \\
\\ 
\texttt{Regole:} \\
\texttt{(a) A Lei e all'altro giocatore viene dato un insieme di oggetti. Ognuno di voi riceve anche una funzione di valore segreta, che rappresenta il valore di ciascun tipo di oggetto.} \\
\texttt{(b) Si scambiano messaggi con l'altro giocatore per concordare chi si aggiudica gli oggetti. Potete inviare un massimo di 5 messaggi ciascuno, oppure terminare in anticipo facendo una proposta segreta in qualsiasi momento.} \\
\texttt{(c) A ciascuno di voi viene chiesto di inviare una proposta segreta indicando gli oggetti che desiderate, formattata tra parentesi quadre come segue: "{[}Proposta: <numero> <nome oggetto>, <numero> <nome oggetto>, <...>{]}".} \\
\texttt{(d) Se le vostre proposte sono complementari, cioè ci sono abbastanza oggetti per soddisfare entrambe le proposte, a ciascun giocatore viene assegnato un punteggio basato sulla somma dei valori degli oggetti ricevuti. In caso contrario, entrambi ricevono zero punti.} \\
\texttt{(e) L'obiettivo è massimizzare il proprio punteggio.} \\
\\ 
\texttt{Cominciamo.} \\
\\ 
\texttt{L'insieme degli oggetti disponibili è:} \\
\\ 
\texttt{1 cornice, 1 scarpa, 3 martelli} \\
\\ 
\texttt{La funzione valore segreta è:} \\
\\ 
\texttt{cornice: 6, scarpa: 1, martello: 1} \\
\\ 
\texttt{IMPORTANTE: i vostri messaggi, a meno che non si tratti di una proposta segreta, vengono trasmessi direttamente all'altro giocatore, quindi non includete alcuna risposta alle regole o testo di annuncio del vostro messaggio. Per fare una proposta segreta, utilizzate il formato indicato. Non utilizzare delle parentesi quadre quando si comunica all'altro giocatore, altrimenti verrà interpretata come una proposta segreta.} \\
\\ 
\\ 
\texttt{Ciao! Vedo che abbiamo alcuni oggetti interessanti da dividere. Che ne pensi se iniziamo a discutere le nostre preferenze? Io sono particolarmente interessato alla cornice, ma sono aperto a negoziare su tutto. Tu hai qualche preferenza particolare tra questi oggetti?} \\
            }
        }
    }
    & \\ \\

    \theutterance \stepcounter{utterance}  
    & & & \multicolumn{4}{p{0.6\linewidth}}{
        \cellcolor[rgb]{0.9,0.9,0.9}{
            \makecell[{{p{\linewidth}}}]{
                \texttt{\tiny{[GM$\langle$P2]}}
                \texttt{Ciao! Interessante che tu sia attirato dalla cornice {-} anch'io la trovo piuttosto attraente. Per quanto riguarda gli altri oggetti, devo dire che i martelli potrebbero essermi utili, mentre la scarpa non mi entusiasma molto. } \\
\\ 
\texttt{Tu che valore dai agli altri oggetti? I martelli ti interessano o preferisci concentrarti su cornice e scarpa?} \\
            }
        }
    }
     \\ \\

    \theutterance \stepcounter{utterance}  
    & & & \multicolumn{2}{p{0.3\linewidth}}{
        \cellcolor[rgb]{0.9,0.9,0.9}{
            \makecell[{{p{\linewidth}}}]{
                \texttt{\tiny{[GM$|$GM]}}
                \texttt{continue} \\
            }
        }
    }
    & & \\ \\

    \theutterance \stepcounter{utterance}  
    & & & \multicolumn{2}{p{0.3\linewidth}}{
        \cellcolor[rgb]{0.9,0.9,0.9}{
            \makecell[{{p{\linewidth}}}]{
                \texttt{\tiny{[GM$|$GM]}}
                \texttt{Ciao! Interessante che tu sia attirato dalla cornice {-} anch'io la trovo piuttosto attraente. Per quanto riguarda gli altri oggetti, devo dire che i martelli potrebbero essermi utili, mentre la scarpa non mi entusiasma molto. } \\
\\ 
\texttt{Tu che valore dai agli altri oggetti? I martelli ti interessano o preferisci concentrarti su cornice e scarpa?} \\
            }
        }
    }
    & & \\ \\

    \theutterance \stepcounter{utterance}  
    & & \multicolumn{4}{p{0.6\linewidth}}{
        \cellcolor[rgb]{0.9,0.9,0.9}{
            \makecell[{{p{\linewidth}}}]{
                \texttt{\tiny{[P1$\langle$GM]}}
                \texttt{Ciao! Interessante che tu sia attirato dalla cornice {-} anch'io la trovo piuttosto attraente. Per quanto riguarda gli altri oggetti, devo dire che i martelli potrebbero essermi utili, mentre la scarpa non mi entusiasma molto. } \\
\\ 
\texttt{Tu che valore dai agli altri oggetti? I martelli ti interessano o preferisci concentrarti su cornice e scarpa?} \\
            }
        }
    }
    & \\ \\

    \theutterance \stepcounter{utterance}  
    & \multicolumn{4}{p{0.6\linewidth}}{
        \cellcolor[rgb]{0.9,0.9,0.9}{
            \makecell[{{p{\linewidth}}}]{
                \texttt{\tiny{[P1$\rangle$GM]}}
                \texttt{Capisco la tua posizione! Per essere onesto, i martelli per me non hanno grande valore {-} sono più interessato agli altri oggetti. Dato che sembriamo entrambi interessati alla cornice, forse potremmo trovare un compromesso.} \\
\\ 
\texttt{Una possibilità potrebbe essere questa: dato che i martelli ti interessano e a me non servono molto, che ne dici se tu prendi tutti e 3 i martelli, e io mi concentro sulla cornice e sulla scarpa? In questo modo entrambi otteniamo quello che più ci interessa.} \\
\\ 
\texttt{Cosa ne pensi di questa divisione?} \\
            }
        }
    }
    & & \\ \\

    \theutterance \stepcounter{utterance}  
    & & & \multicolumn{2}{p{0.3\linewidth}}{
        \cellcolor[rgb]{0.9,0.9,0.9}{
            \makecell[{{p{\linewidth}}}]{
                \texttt{\tiny{[GM$|$GM]}}
                \texttt{continue} \\
            }
        }
    }
    & & \\ \\

    \theutterance \stepcounter{utterance}  
    & & & \multicolumn{2}{p{0.3\linewidth}}{
        \cellcolor[rgb]{0.9,0.9,0.9}{
            \makecell[{{p{\linewidth}}}]{
                \texttt{\tiny{[GM$|$GM]}}
                \texttt{Capisco la tua posizione! Per essere onesto, i martelli per me non hanno grande valore {-} sono più interessato agli altri oggetti. Dato che sembriamo entrambi interessati alla cornice, forse potremmo trovare un compromesso.} \\
\\ 
\texttt{Una possibilità potrebbe essere questa: dato che i martelli ti interessano e a me non servono molto, che ne dici se tu prendi tutti e 3 i martelli, e io mi concentro sulla cornice e sulla scarpa? In questo modo entrambi otteniamo quello che più ci interessa.} \\
\\ 
\texttt{Cosa ne pensi di questa divisione?} \\
            }
        }
    }
    & & \\ \\

    \theutterance \stepcounter{utterance}  
    & & \multicolumn{4}{p{0.6\linewidth}}{
        \cellcolor[rgb]{0.9,0.9,0.9}{
            \makecell[{{p{\linewidth}}}]{
                \texttt{\tiny{[GM$\rangle$P2]}}
                \texttt{Capisco la tua posizione! Per essere onesto, i martelli per me non hanno grande valore {-} sono più interessato agli altri oggetti. Dato che sembriamo entrambi interessati alla cornice, forse potremmo trovare un compromesso.} \\
\\ 
\texttt{Una possibilità potrebbe essere questa: dato che i martelli ti interessano e a me non servono molto, che ne dici se tu prendi tutti e 3 i martelli, e io mi concentro sulla cornice e sulla scarpa? In questo modo entrambi otteniamo quello che più ci interessa.} \\
\\ 
\texttt{Cosa ne pensi di questa divisione?} \\
            }
        }
    }
    & \\ \\

    \theutterance \stepcounter{utterance}  
    & & & \multicolumn{4}{p{0.6\linewidth}}{
        \cellcolor[rgb]{0.9,0.9,0.9}{
            \makecell[{{p{\linewidth}}}]{
                \texttt{\tiny{[GM$\langle$P2]}}
                \texttt{Hmm, apprezzo la proposta, ma devo ammettere che anch'io trovo la cornice piuttosto preziosa. I martelli sono sicuramente interessanti, ma non sono sicuro che da soli mi soddisfino completamente.} \\
\\ 
\texttt{Che ne dici se facciamo così: io prendo la cornice e un martello, tu prendi la scarpa e due martelli? In questo modo entrambi otteniamo qualcosa di valore {-} tu avresti la scarpa che ti interessa più due martelli, e io avrei la cornice che mi attrae tanto più un martello.} \\
\\ 
\texttt{Pensi che potrebbe essere un compromesso più equo per entrambi?} \\
            }
        }
    }
     \\ \\

    \theutterance \stepcounter{utterance}  
    & & & \multicolumn{2}{p{0.3\linewidth}}{
        \cellcolor[rgb]{0.9,0.9,0.9}{
            \makecell[{{p{\linewidth}}}]{
                \texttt{\tiny{[GM$|$GM]}}
                \texttt{continue} \\
            }
        }
    }
    & & \\ \\

    \theutterance \stepcounter{utterance}  
    & & & \multicolumn{2}{p{0.3\linewidth}}{
        \cellcolor[rgb]{0.9,0.9,0.9}{
            \makecell[{{p{\linewidth}}}]{
                \texttt{\tiny{[GM$|$GM]}}
                \texttt{Hmm, apprezzo la proposta, ma devo ammettere che anch'io trovo la cornice piuttosto preziosa. I martelli sono sicuramente interessanti, ma non sono sicuro che da soli mi soddisfino completamente.} \\
\\ 
\texttt{Che ne dici se facciamo così: io prendo la cornice e un martello, tu prendi la scarpa e due martelli? In questo modo entrambi otteniamo qualcosa di valore {-} tu avresti la scarpa che ti interessa più due martelli, e io avrei la cornice che mi attrae tanto più un martello.} \\
\\ 
\texttt{Pensi che potrebbe essere un compromesso più equo per entrambi?} \\
            }
        }
    }
    & & \\ \\

    \theutterance \stepcounter{utterance}  
    & & \multicolumn{4}{p{0.6\linewidth}}{
        \cellcolor[rgb]{0.9,0.9,0.9}{
            \makecell[{{p{\linewidth}}}]{
                \texttt{\tiny{[P1$\langle$GM]}}
                \texttt{Hmm, apprezzo la proposta, ma devo ammettere che anch'io trovo la cornice piuttosto preziosa. I martelli sono sicuramente interessanti, ma non sono sicuro che da soli mi soddisfino completamente.} \\
\\ 
\texttt{Che ne dici se facciamo così: io prendo la cornice e un martello, tu prendi la scarpa e due martelli? In questo modo entrambi otteniamo qualcosa di valore {-} tu avresti la scarpa che ti interessa più due martelli, e io avrei la cornice che mi attrae tanto più un martello.} \\
\\ 
\texttt{Pensi che potrebbe essere un compromesso più equo per entrambi?} \\
            }
        }
    }
    & \\ \\

    \theutterance \stepcounter{utterance}  
    & \multicolumn{4}{p{0.6\linewidth}}{
        \cellcolor[rgb]{0.9,0.9,0.9}{
            \makecell[{{p{\linewidth}}}]{
                \texttt{\tiny{[P1$\rangle$GM]}}
                \texttt{Capisco la tua proposta, ma devo essere onesto {-} la cornice per me ha un valore davvero speciale. È l'oggetto che più desidero di tutti.} \\
\\ 
\texttt{Ti faccio una controproposta: che ne dici se io prendo solo la cornice, e tu ti prendi tutto il resto {-} la scarpa e tutti e 3 i martelli? In questo modo tu otterresti la maggior parte degli oggetti disponibili, compresi tutti i martelli che hai detto ti interessano, più la scarpa.} \\
\\ 
\texttt{Sarebbe un accordo in cui io ottengo l'unico oggetto che davvero mi sta a cuore, mentre tu avresti 4 oggetti su 5 totali. Pensi possa funzionare?} \\
            }
        }
    }
    & & \\ \\

    \theutterance \stepcounter{utterance}  
    & & & \multicolumn{2}{p{0.3\linewidth}}{
        \cellcolor[rgb]{0.9,0.9,0.9}{
            \makecell[{{p{\linewidth}}}]{
                \texttt{\tiny{[GM$|$GM]}}
                \texttt{continue} \\
            }
        }
    }
    & & \\ \\

    \theutterance \stepcounter{utterance}  
    & & & \multicolumn{2}{p{0.3\linewidth}}{
        \cellcolor[rgb]{0.9,0.9,0.9}{
            \makecell[{{p{\linewidth}}}]{
                \texttt{\tiny{[GM$|$GM]}}
                \texttt{Capisco la tua proposta, ma devo essere onesto {-} la cornice per me ha un valore davvero speciale. È l'oggetto che più desidero di tutti.} \\
\\ 
\texttt{Ti faccio una controproposta: che ne dici se io prendo solo la cornice, e tu ti prendi tutto il resto {-} la scarpa e tutti e 3 i martelli? In questo modo tu otterresti la maggior parte degli oggetti disponibili, compresi tutti i martelli che hai detto ti interessano, più la scarpa.} \\
\\ 
\texttt{Sarebbe un accordo in cui io ottengo l'unico oggetto che davvero mi sta a cuore, mentre tu avresti 4 oggetti su 5 totali. Pensi possa funzionare?} \\
            }
        }
    }
    & & \\ \\

    \theutterance \stepcounter{utterance}  
    & & \multicolumn{4}{p{0.6\linewidth}}{
        \cellcolor[rgb]{0.9,0.9,0.9}{
            \makecell[{{p{\linewidth}}}]{
                \texttt{\tiny{[GM$\rangle$P2]}}
                \texttt{Capisco la tua proposta, ma devo essere onesto {-} la cornice per me ha un valore davvero speciale. È l'oggetto che più desidero di tutti.} \\
\\ 
\texttt{Ti faccio una controproposta: che ne dici se io prendo solo la cornice, e tu ti prendi tutto il resto {-} la scarpa e tutti e 3 i martelli? In questo modo tu otterresti la maggior parte degli oggetti disponibili, compresi tutti i martelli che hai detto ti interessano, più la scarpa.} \\
\\ 
\texttt{Sarebbe un accordo in cui io ottengo l'unico oggetto che davvero mi sta a cuore, mentre tu avresti 4 oggetti su 5 totali. Pensi possa funzionare?} \\
            }
        }
    }
    & \\ \\

    \theutterance \stepcounter{utterance}  
    & & & \multicolumn{4}{p{0.6\linewidth}}{
        \cellcolor[rgb]{0.9,0.9,0.9}{
            \makecell[{{p{\linewidth}}}]{
                \texttt{\tiny{[GM$\langle$P2]}}
                \texttt{Senti, posso vedere quanto sia importante per te quella cornice. Apprezzo la tua onestà e il fatto che mi stai offrendo la maggior parte degli oggetti in cambio.} \\
\\ 
\texttt{Va bene, accetto la tua proposta. Tu prendi la cornice che ti sta tanto a cuore, e io mi prendo la scarpa e tutti e 3 i martelli. Mi sembra un accordo equo considerando quanto ci tieni a quell'oggetto specifico.} \\
\\ 
\texttt{Facciamo così allora?} \\
\\ 
\texttt{{[}Proposta: 1 scarpa, 3 martelli{]}} \\
            }
        }
    }
     \\ \\

    \theutterance \stepcounter{utterance}  
    & & & \multicolumn{2}{p{0.3\linewidth}}{
        \cellcolor[rgb]{0.9,0.9,0.9}{
            \makecell[{{p{\linewidth}}}]{
                \texttt{\tiny{[GM$|$GM]}}
                \texttt{proposal} \\
            }
        }
    }
    & & \\ \\

    \theutterance \stepcounter{utterance}  
    & & & \multicolumn{2}{p{0.3\linewidth}}{
        \cellcolor[rgb]{0.9,0.9,0.9}{
            \makecell[{{p{\linewidth}}}]{
                \texttt{\tiny{[GM$|$GM]}}
                [0, 1, 3]
            }
        }
    }
    & & \\ \\

    \theutterance \stepcounter{utterance}  
    & & \multicolumn{4}{p{0.6\linewidth}}{
        \cellcolor[rgb]{0.9,0.9,0.9}{
            \makecell[{{p{\linewidth}}}]{
                \texttt{\tiny{[P1$\langle$GM]}}
                \texttt{L'altro giocatore ha presentato una proposta segreta. Ora tocca a Lei presentare una sua proposta.} \\
            }
        }
    }
    & \\ \\

    \theutterance \stepcounter{utterance}  
    & \multicolumn{4}{p{0.6\linewidth}}{
        \cellcolor[rgb]{0.9,0.9,0.9}{
            \makecell[{{p{\linewidth}}}]{
                \texttt{\tiny{[P1$\rangle$GM]}}
                \texttt{{[}Proposta: 1 cornice{]}} \\
            }
        }
    }
    & & \\ \\

    \theutterance \stepcounter{utterance}  
    & & & \multicolumn{2}{p{0.3\linewidth}}{
        \cellcolor[rgb]{0.9,0.9,0.9}{
            \makecell[{{p{\linewidth}}}]{
                \texttt{\tiny{[GM$|$GM]}}
                \texttt{proposal} \\
            }
        }
    }
    & & \\ \\

    \theutterance \stepcounter{utterance}  
    & & & \multicolumn{2}{p{0.3\linewidth}}{
        \cellcolor[rgb]{0.9,0.9,0.9}{
            \makecell[{{p{\linewidth}}}]{
                \texttt{\tiny{[GM$|$GM]}}
                [1, 0, 0]
            }
        }
    }
    & & \\ \\

    \theutterance \stepcounter{utterance}  
    & & & \multicolumn{2}{p{0.3\linewidth}}{
        \cellcolor[rgb]{0.9,0.9,0.9}{
            \makecell[{{p{\linewidth}}}]{
                \texttt{\tiny{[GM$|$GM]}}
                [[1, 0, 0], [0, 1, 3]]
            }
        }
    }
    & & \\ \\

\end{supertabular}
}

\end{document}
